% \subsection{Implementation}
% \label{sec:implementation}


% \subsubsection{Algorithm Overview}
% \label{sec:algorithm}

% The resulting algorithm as described by \textcite{deliot2023} that integrates all the ideas and insights from the previous chapters works as follows:
% %
% \todo{Rewrite}
% \begin{enumerate}
%     \item First calculate the footprint of the pixel and parametrize it into the three dimensions: phi (orientation), anisotropy, and level of detail.
%     \item Find the heptahedron in 3D footprint space that contains the current pixel's footprint.
%     \item Find the tetrahedron in this heptahedron that contains the current pixel's footprint and calculate the barycentric coordinates.
%     \item Sample each of the 4 vertices of the tetrahedron and interpolate the result using the Distributed Binomial Law. Each vertex is sampled using the following steps:
%           \begin{enumerate}
%               \item Compensate the uv parametrization of the vertex for its orientation, anisotropy, and level of detail.
%               \item In the compensated uv space, find the uv triangle that contains the compensated uv coordinates.
%               \item For each triangle vertex and half-vector draw a binomial sample.
%           \end{enumerate}
% \end{enumerate}

% % TODO: Figure showing pipeline geometrically and/or visually

% \subsubsection{Implementation Details}
% \label{sec:implementation-details}